% ================================================================
%  第 10 章 孩子嬉戏 (原书 p.209-218, 边码 111-121)
%  源文: _work/src1/task-ch10.md
% ================================================================
% ---------------- p.209 ----------------
\chapter{孩子嬉戏}

\section{(42) 孩子}

肯定还有些角落,扫帚没扫到。没关系,它们自会引起我的注意,到那时再收拾也不迟。至于我那大名鼎鼎的“数学家过往”,大扫除已经做完,这一点毫无疑问。

既然我又一次看清我并不比旁人更好,那就别再掉进那个老掉牙的陷阱:把自己看成\textbf{比我自己更好}!——把如今的我,退出了那台旋转木马、诸如此类的我,看成比十五年前、或者十五天前的我更好。这十五年里我学到了些东西,这是肯定的;这十五天里也是,甚至从昨天起也是。我学到东西的时候,就成熟了,就不再完全是同一个人了。学到某样东西之后的我,并不比那样东西还摆在我面前时的我“更好”。一个更熟的果子并不比一个不太熟的、或者青的果子“更好”。一个季节并不比它前一个季节“更好”。最熟的果子味道可能更可口,也可能更不可口,这要看各人的口味。一年一年过去,我在自己的皮囊里觉得更自在了,想必我身上发生的变化“合我的口味”——可这些变化并不合我所有朋友或亲近者的口味。每次我重新做起数学,四面八方都来恭维,腔调是这样的:“他当初干吗要去弄别的东西呢!现在一切都归了位,早该如此!”看着一个人变,总叫人不安……

\origpage{111} 我学习,我成熟,我变化——以至于有时我很难在从前那个人身上认出自己,得靠一段回忆、或者旁人一次意想不到的见证,才重新发现他。我在变;同时也有某种东西始终“是同一个”。它一直都在,从永远起就在,肯定从我出生起,也许还更早。近几年,我觉得自己渐渐能把它认准了。我把它叫作“孩子”。凭这东西,此刻的我不比一生中任何别的时刻更好;它就在那里,尽管往往很难猜出它的在场。也凭这东西,\textbf{我不比任何人更好},也没有\textbf{任何人比我更好}。在某些时刻、在某些人身上,孩子更在场一些。而这是件很让人受用的事。这并不意味着某个人比另一个人“更好”,或者比他自己的另一个时刻“更好”。

常常,我做数学的时候,或者做爱的时候,或者冥想的时候,是孩子在玩。并不总是只有他一个在“玩”。可他不在的时候,就没有数学,没有爱,也没有冥想。装样子是没用的——这出戏我很少演。

% ---------------- p.210 ----------------
\origpage{112} 当然不只有孩子。还有“我”,还有“老板”,或者说“大老板”,爱怎么叫就怎么叫。对企业运转来说,老板肯定不可或缺。既然有个老板,总该是为着什么。他照管庶务,而且像所有老板一样,有个讨厌的倾向:越占越宽。他把自己看得无比要紧,非要处处压过对面的老板。占得宽也好,不宽也好,他不过是老板,他不是工人。他组织,他指挥,他也收账——当然!收益他当作自己应得之物收下,亏损他当作侮辱来承受。但他什么也不创造。只有工人有创造的力量,而工人不是别人,正是孩子。

老板和工人合得来的企业,很少见。多数时候,根本看不到工人的影子,他被关在天知道什么地方。是老板做出顶替他站在车间里的样子,结果可想而知。也常常是这样:工人明明就在,老板却对他开战,或是猛烈的战争,或是小股交火——这样的车间里出不了什么像样的东西!有时老板对工人也有一种带戒心的宽容,他一边嘀咕一边由着他干,眼睛却一刻也不离开他。那就像一场从未停歇的战争里,不断续签的停战。工人也就趁着这停战,勉强干上一点活。

我刚才那场冥想的功效,绝不意味着我对数学的那份占有态度就此如中了魔法一般消失了!至少,我得近得多地去看占有欲的种种表现——刚才我不过是叫出了其中一种的名字,轻轻碰到了它一下。这里不是地方:这篇“导论”已经变成了“导论性的一章”,而这一章自己也已经显得长了!不过昨夜有一件事“咔嗒”一声对上了,此刻我想稍稍回过头去说它一说;那是两三年前我带着几分惊讶记下的一件事。

当时我正钻在一个数学问题上,已说不上是什么问题了;某一刻(不知由于什么情况),我发现我盯着的问题也许早有人盯过,它很可能在某本书里白纸黑字地处理过,只消我去图书馆查一查。单是想到这种可能性,效果就如雷击,把我惊得发呆:转眼之间,欲望不见了。忽然之间,那个我也许已经花了好几星期、还打算再花上几星期的问题,对我完全失去了吸引力!那不是赌气,那是兴趣突然而彻底的丧失。就算那本书就在我手里,我也懒得翻开。

\origpage{113} 事实上,那种可能性并未得到证实,于是欲望又回来了,我照旧一路做下去,像什么都没发生过。我却仍然愣着。当然,如果我正在做的东西真真是我另做一件事所\textbf{必需}的,那就不会有那么惊人的兴趣跌落。我常重做已知的东西,明知或猜到它们已知,却一点也不在乎。那时我正处在一种冲劲里:照我自己的方式、按它们向我呈现的样子去做,比到书里、文章里去翻找更省事,尤其有意思得多。那时我是“顺势”奔向别的东西,奔向欲望把我带往的地方。当然,我够“在行”,知道自己最终要到的那个地方,任何书、任何文章里都没有。

这让我想起来:数学工作即使孤零零地做了许多年,也不是纯粹个人的、个体性的工作,不像冥想那样——至少在我身上不是。我在数学里追的那个“未知”,要能那样强有力地吸引我,就必须不只对我未知,而是对所有人都未知。数学书里写着的东西不是未知,即便我自己从未听说过它。读书或读文章从来不吸引我,只要\textbf{能躲我就躲}。它能告诉我的,从来都不是未知;
% ---------------- p.211 ----------------
我对它的兴趣,也没有欲望的那种质地。那是一种“权宜的”兴趣,是对可能对我有用的一条\textbf{信息}的兴趣,是某个欲望的工具,而这条信息根本不是那个欲望的对象。

细想之后,我不觉得我讲的那件事是什么嫉妒、占有的心性的迹象,也不是什么受了挫的虚荣的迹象。我心里没有半点赌气,没有半点失望,只是一个前一瞬还强烈的欲望突然消失了。那是在一个我压根没想过要发表任何东西的时期,也没想过有一天我会忽然又想发表点什么。这欲望不是虚荣的表现,不是急于积累知识、头衔和信誉的饕餮——它确确实实是一个真实的欲望,是沉迷于游戏的孩子的那种欲望。而忽然之间——什么都没有了!谁懂就懂吧,我是不懂……抱歉!

\section{(43) 扫兴的老板——或高压锅}

\origpage{114} \noindent\pmark 我觉得,这场对我数学家一生的回顾,终究是做完了。当然,我没有把题目做尽——那得要好几卷书,假如这样的题目真能“做尽”的话。那不是我的意图。我的意图是弄清楚:在某种我如今只能一阵一阵感觉到的“气息”的出现里,我到底是不是参与者、是不是共同行动者;如果是,又是以什么方式。现在我弄清楚了,这让人舒坦。再往前走,去深挖那些只是瞥见或触及的东西,可能会很来劲。可看、可做、可发现的有意思的东西太多了!至于我作为数学家的过去,我觉得,为了承担这段过去而\textbf{必须}看的,已经看到了。

肯定,若把这番冥想往深里做,我准会学到许多关于我当下的有意思的东西。这项工作几乎每走一步都让我感到的一件事是:我对这段过去还眷恋得多么深,它直到今天在我的自我形象里、以及在我与他人的关系里,分量有多重;尤其是在我与那些——从某种意义上说——被我抛下的人的关系里。这番工作过程中,我与这段过去的关系肯定变了,变得疏离了些,或者说轻了些。往后自有分晓。但很可能还会留有一份眷恋,只要我的数学激情还没烧尽、还没餍足——只要我还“做数学”。至于它会不会先我而熄灭,我一点也不想去猜、去预言……

十多年里,我一直以为这激情已经熄灭。更确切地说,是我径自宣布了它已熄灭。那正是我暂时停下做数学、并且重新发现世界的那一天!此后三四年,我全力投入一项如此紧张的活动,旧日的激情想必找不到一丝缝隙钻出来现身。那是紧张学习的几年,不过在某个层面上还相当浮浅。在那之后的几年里,数学激情以突如其来的、完全无法预料的发作显现出来。每次发作持续几个星期或几个月,而我执意不去理会它们那相当清楚的含义。我早就一次性决定了:做数学的那股馋劲,既然明摆着百无一用,从此就是过时的玩意儿,句号!可那“百无一用”的东西并不这么听——而我这边,照旧充耳不闻。

\origpage{115} \noindent\pmark 有件事也许显得悖谬:恰恰是在发现冥想之后(1976年)、在我的生命里进来一种新的激情之后,旧激情才重新露面,而且格外猛烈,几乎算得上暴烈——仿佛每次都是一只锅盖在过强的压力下崩开。又过了五年,在形势的推动之下——这话说得不假——我才费心去察看究竟发生了什么。那是我做过的最长的一次冥想,针对一个看来范围分明的问题:我花了六个月顽强而紧张的工作,才绕着一座冰山走了一圈;它那可见的尖顶终于碍事到逼着我——几乎是违心——去看个究竟。不得不确认的是一种冲突局面,而从表面看,那是两股力量、两种欲望的冲突:冥想的欲望,和做数学的欲望。

% ---------------- p.212 ----------------
在这场漫长的冥想中,我一步一步学到:做数学的欲望——我一向鄙弃它——和冥想的欲望一样,都是孩子的欲望,而后者我捧到了天上。孩子才不管大老板、大掌柜那份鄙弃,也不管他那份谦逊的傲气!孩子的欲望一个接一个,随着时辰和日子,像一支舞的动作,一个从另一个里生出来。它们的本性就是如此。它们并不彼此对立,正如一首歌的各个诗节并不彼此对立,或者一首康塔塔、一首赋格的各个乐章并不彼此对立。是老板这个蹩脚的乐队指挥,宣布这个乐章“好”、那个乐章“坏”,在本来和谐的地方制造出冲突。

这场冥想之后,老板安分了些,不太做出把鼻子伸到不关他事之处的样子。这一回工作拖得很长,而我原以为几天就能完。工作一做完,“结果”就显得明摆着,几句话就能说出来\fnoteP{1}{(37)}{我想,几乎不必补充:这项旷日持久的工作,一天一天地还显露出别的东西,远不止我刚才以碑铭式的简短形式交出的那个“结果”。冥想的工作如此,为一个打算考察的特定问题所推动的数学工作也莫不如此。所走的那条路(它通向、或不通向对原问题或多或少的澄清)上的种种曲折,往往比原问题本身、或比“最终结果”更有意思。} (37)。可要是有个眼尖的人,在工作之前或进行中就把这几句话告诉我,恐怕也一点都帮不上我。工作之所以拖这么久,是因为阻力又强又深。再说,老板这回被骂得狗血淋头,却一声没吭,因为当时的气氛根本不容他发火。可以肯定的是,这六个月用得其所,而且是省不掉的;正如一个女人省不掉九个月的妊娠,到头来生下一个像娃娃一样“明摆着”的东西。

\section{(44) 再次倒转蒸汽}

到这里,我已经一年半没冥想了,只在十二月里有那么几个小时,为的是看清一个紧迫的问题。而这一年来,我把大部分精力投在数学上。这一“潮”来得和别的潮一样,数学潮也好,冥想潮也好:它们来时不打招呼。就算打了招呼,我也从来听不见!老板对冥想还留着一点小小的偏好,看来不得不信:每一次,我原以为会永远持续的冥想潮,转眼已被数学潮接替;而那数学潮——在我看来不过是几天、最多几星期的事——却拖宕下去,铺展成几个月,也许甚至,谁知道呢,成几年。不过老板终于明白:这些节律不是他造出来的,想去调节它们,他什么也捞不着。

但也许老板那点“小小的偏好”终究翻转过来了,因为近一年来事情已经说定、已经决定:我至少要用几年工夫去“重做数学”,可以说是正式地——我甚至申请了国家科研中心(CNRS)的一个职位!更重要、而且一年前还完全想不到的是:我又开始发表东西了。即便在刚才提到的1981年那场冥想之后,在做数学的欲望不再被当作穷亲戚对待之后,我也不会想到自己还能重新发表数学。顶多是别的东西:一本谈冥想的书,或者谈梦与做梦者的书——再说,我当时忙得根本没闲心去写一本关于它的书!
% ---------------- p.213 ----------------
写它干吗?!

于是就有了一个相当重要的决定,它牵动我今后若干年的生命进程,而它是在一种捎带着的情形下作出的,我甚至说不清是何时、怎样作出的。\origpage{116}\origpage{117} 有一天,当打字的笔记已经攒成厚厚一沓时(瞧瞧!直到那时,我那些数学上的苦思,我都还只是手写了事……\fnoteP{2}{(38)}{这些笔记其实是那封写给……的长信的继续,那封信成了它的第一章。它们被打字机打出来,是为了让这位旧日的朋友、以及另外两三位(尤其是龙尼·布朗(Ronnie Brown))读得清楚,我本以为他们可能会感兴趣。再说,那封信从未收到回音,收信人根本没有读它;将近一年之后(我问他是否收到了),他竟真心感到惊讶:我居然会有一刻以为他会读它,鉴于人们该期待我做的那种数学……} (38),关于栈、同伦模型等等),事情就成了定局:这个,发表!既然要做,索性做足,开一个小系列的数学反思,名字现成,只要把字母大写:“数学反思”!大致就是这样——此刻那团大名鼎鼎的“迷雾”向我复述的,就是这些;这迷雾常常替我充当记忆。这一回的回忆,肯定大大删削过了。无论如何,值得注意的是:这件事做成的时候,我连停下来\textbf{看看}自己往哪儿走、是什么在推我、带我的工夫都没有……这倒是我还想再做一次的:顺着这场不期而至的冥想的势头,好让我能觉得它真正做完了。

立刻浮上心头的问题是:我刚指出的这件“值得注意的事”,是老板那种(所谓的?)“谨慎”的迹象——他无论如何都不肯插手(哪怕只是投去一瞥不检点的目光……)一个如此美妙、根本不需要他的自发运动,如此等等;还是\textbf{相反},这迹象表明他已经断然站了边,而那所谓的“小小的偏好”正把他朝数学的方向全力推?

只消把问题白纸黑字写下来,答案就自己出现了!并不是那小子——他这一回玩起的把戏也许比别的把戏更长久——就此宣布自己要连着 X 年不停手地干下去,老老实实地涂黑纸张,涂够所需的那么多页,凑成漂亮一套大写字头丛书里数目可观的若干卷!是老板把一切都预见到、安排好了,小子只消照办。也许小子巴不得如此,这没法预先知道——但那是次要问题。再说,小子的兴致至少在一定程度上取决于\textbf{境况},而境况主要取决于老板。

老板选了边,再清楚不过。再说他刚显出过某种灵活:一个多月来,一场冥想就在他和蔼的目光下进行。诚然,他的和蔼绝非无私,因为冥想的有形成果,我正在写的这些笔记,将成为他已经看见自己建起的那座塔最美的奠基石,而砌塔的石头,是那位看来很听话的工人-孩子乖乖凿好的。说真的,现在夸他“灵活”还早了点儿!三个月前那几小时冥想,一年半里头总共就这么点,未免太寒碜了!

\origpage{118} \noindent\pmark 然而,我并不觉得这一整段时间里,有过什么被压抑、被挫败的冥想欲望。十二月的那几个小时里,我把事情理清了,看见了该看见的;这就足以改变一个本来不清楚的局面。我重新接上中断的数学工作,而不必把别的什么一刀斩断。我不觉得有冲突蹑手蹑脚地又回来了——我是说,那个两年多前已经化解、这一次会以颠倒的形式重新出现的冲突。老板有偏好,这是他本性如此,也完全是他权利——他要是装作禁止自己有偏好,那才蠢呢(虽然比这更蠢的事也有……)。这不是冲突的迹象,尽管冲突常由它而起。事情到了这一步,实在看不出有什么该怪他不灵活的!

% ---------------- p.214 ----------------
这一点看清了,剩下的是试着圈定老板的“动机”——为了这次蒸汽倒转;它做得再低调不过,可仔细看来,其实相当惊人。

\section{(45) 不是上师的上师——或三条腿的马}

这立刻把我带回1981年7月到12月的那场冥想;在那之前,我刚在一种数学狂热里过了四个月。这段有点疯的日子(从数学上看还颇有收获\fnoteP{3}{(39)}{这尤其是《穿越伽罗瓦理论的长征》(La Longue Marche à travers la théorie de Galois)的时期,《纲领概要》(Esquisse d'un Programme)里谈到过它(第3节:“与儿童画相关联的数域”)。} (39))随着一个梦,一夜之间结束了。那梦以一个具有野性而不可抗拒之力的比喻,描绘了我生命里正在发生的事——一个关于这场狂热的比喻。信息明如闪电,可我还是花了两天紧张的工作,才接受它那明摆着的含义\fnoteP{4}{(40)}{\emph{来访}\par 关于这个梦的工作,是一封用英文写的长信,写给一位朋友兼同行,他前一天曾在我家匆匆过了一趟。做梦者用来从一片表面的虚无里唤起这个逼真得惊人的梦的某些材料,显然取自这位亲爱的朋友那次短暂来访——我已经快十年没见过他了。因此,工作的第一天,与我以往的经验相反,我竟以为可以断定:我做的这个梦关系到我这位朋友,多于关系到我自己——该做这个梦的是他,不是我!这是一种回避梦境信息的方式,而那信息(我本该凭过去的经验一下子就明白)只关乎我一个人。我最终在工作这第一阶段——浮浅的阶段——之后的夜里明白了这一点;第二天我在同一封信里接着写下去。自那封值得记住的信之后,我再没有收到过这位朋友的任何音讯,他是我有过的朋友中最亲近者之一。\par 这次工作也是唯一一次以信件形式(而且还用的是英文)进行的冥想,因此我手上再没有它的书面痕迹。这个插曲让我格外震动——在许多别的插曲之中,它们都表明:一切超出某种门面的工作的迹象,一切把极简单、却通常被大家当作义务去忽视的事实摆到光天化日之下的工作,会在旁人心里引起多大的不适与恐惧。这一点我后面还要回来谈(见第47节,“孤独的历险”)。} (40)。这之后,我知道了自己该做什么。此后六个月的工作里,我再没回到那个梦上去,可我所做的,无非是更深入地进入它的含义、完整地消化它的信息。梦后第三天,这信息只在相当表面、相当粗糙的层面上被理解了。我尤其需要深挖的,是“我的”关系——我是说老板的——与当下这两种欲望各自的关系,而这两种欲望在我看来相互敌对。

\origpage{119} 自那场冥想以来,我生命里发生了这么多事,以至于它显得像是属于一个很遥远的过去。如果我试着把从它那里学到的、关于“老板”动机的东西说出来,就是这样:从“第一次觉醒”(1970年)算起、到那时已经过去的十二年里,老板押的注,明摆着是“\textbf{那匹劣马}”:\textbf{在数学与冥想之间}(他喜欢把两者对立起来)\textbf{他选了冥想}。

 这是一种说法,因为“冥想”这东西连同这个名字,只是五年前的1976年10月才进入我的生命。但在1970年那副刷新过的、我珍爱的自画像里,冥想六年之后来得正是时候,用它的光彩烘托出某种态度或姿态——这姿态早已被察觉,却直到1981年这场冥想才被审视。我把它叫作“大师综合征”,也有人(有道理地)称之为我的“上师做派”。我采用前一个称呼而不是后一个,多半是因为它便于把事情的性质搅混,而我乐意让自己耽在这种混淆里。我身上确实从小就有一个教人的自发乐趣,它与学的自发乐趣毫不相悖;
% ---------------- p.215 ----------------
也毫无做作的成分。在我与学生的关系里起作用的,主要就是这股力量;这关系是浮浅的,却强而纯正——我所谓“纯正”,是指没有做作。是在我称为1970年“觉醒”之后——一个我曾熟悉的天地退缩到几乎消失,连同学生、连同那些让我“教”点什么、让我把我知道的、对我有意义有价值的东西讲出来的机会一起退去——正是在那时,“老板”尽其所能地报了仇:他不再去教数学(那玩意儿只配用来糊口,除此之外配不上我新的伟大),我仿佛看见自己在用我的生活与榜样,教某种“智慧”。当然,我很小心,无论对自己还是对别人都不把这话说出来;而当我听到这个方向的回声时,我当然要矢口否认,并为这些朋友或亲近者如此不理解我而难过。我怎么解释也没用,他们就是不肯懂,真是些再叫人遗憾不过的学生!

\origpage{120} 我读过克里希那穆提(Krishnamurti)的一两本书,印象很深,脑子转眼就把某种信息和某些价值吸收了下来\fnoteP{5}{(41)}{\emph{克里希那穆提,或变成了羁绊的解放}\par 要说我从这次阅读里得到的只是一套词汇,以及一种把它据为己有、最终又理所当然地用它替换现实的倾向,那并不准确。我拿到手的第一本克里希那穆提的书之所以让我如此震动(何况我不过有闲读了几章),是因为他所说的,彻底掀翻了许多对我而言不言自明的东西;而我立刻发觉,那些都是\textbf{老生常谈},从来就是我呼吸的空气的一部分。同时,这次阅读第一次让我注意到一些影响深远的事实,尤其是“对\textbf{现实}的逃避”这一事实——它是心智最强大、最普遍的制约之一。这给了我一把钥匙,去理解此前一直无法理解、并因此(在五六年之后发现冥想之前,我并未察觉这一点)滋生焦虑的种种处境。我立刻在身边各处看到了这种逃避的实情。这解开了一些焦虑,却没有改变任何本质的东西,因为我只在旁人身上看见那实情,同时(理所当然地)设想它在我身上并不存在,总之我是那验证规则的例外(而且对这个确实了不起的例外,我再没多问一句)。事实上,我对别人、对自己都毫无好奇。这把“钥匙”\textbf{只}在怀着深入之欲的人手里才开锁。在我手里,它变成了驱邪术和做派。\par 是在1974年初,我第一次面对这样一个明证:我生命里的破坏一步步跟着我,不可能只来自别人;我身上有某种东西在吸引它、滋养它、让它延续下去。那是一个谦卑而敞开的时刻,宜于更新。可那次更新当时还停在边缘,而且短暂,因为缺少一种\textbf{深入的工作}。那“\textbf{我身上的某种东西}”还很模糊。我分明看出那是爱的匮乏,但连“做一番工作、更切近地查明我身上的爱在何处、如何匮乏,它如何表现出来,造成了哪些具体后果”等等的念头,都不可能来自我至此认识的任何圈子、任何人,也不可能来自克里希那穆提。(恰恰相反,克氏喜欢强调一切工作的徒劳,他自动把它等同于自我的“急于成为”。)于是,我拿一份借来的“智慧”当全部罗盘,除了耐心等待“爱”像圣灵的恩典一样降临到我身上,我看不出还能做什么。\par 然而,我刚刚在一个浪头的最深处学到的那点朴素的真理,却激起了一个强劲的新能量之浪,其势头可比两年半后推动我第一次投入冥想的那个浪。这股能量当时并没有完全闲置。几个月后,就在我因一场及时降临的事故动弹不得的那时,它推动了一次(写下来的)反思:一生中第一次,我审视了那个曾是我与他人关系之未言明基础的世界观,它来自我的父母,尤其来自我母亲。那时我十分清楚地意识到,这个世界观已经破产,它无法说明人与人之间关系的实情,也无助于我这个人以及我与他人的关系之舒展。这次反思仍带着“克里希那穆提风格”的印记,也带着克氏那条禁忌:任何朝向理解的\textbf{真正的工作}都不许做。可它还是使几个月前诞生、起初模糊而难以捉摸的一种认识变得切实而不可逆转。这种认识,当时世上任何书、任何人都无法给我。\par 这次反思要够得上冥想的品质,最缺的是对我的这个人、对我在我自己眼中的\textbf{形象}投去的目光,而不只是对世界观的目光——也就是对一个公理系统、而我并未真正“有血有肉”地出现在其中的目光。它还缺的,是在\textbf{当下}、在反思进行的那一刻对我自己的目光(这次反思仍停留在真正的工作之下);这目光本会让我同时看出一种借来的文风,以及对这些笔记文学外表某种沾沾自喜,也就是缺乏自发性、缺乏真实。无论它多么不足,对我与他人的关系当下影响也相对有限,这次反思在我看来仍是一个阶段——鉴于起点,这大概是必经的阶段——通向两年后应当发生的更深的更新。正是在那时,我终于发现了冥想——在发现这第一个意想不到的事实时:\textbf{关于我自己这个人,竟还有东西可发现}——那些东西几乎完全决定了我一生的走向,以及我与他人的关系的性质……} (41)。无需再多,就足以让我相信一切都已
% ---------------- p.216 ----------------
到手(当然,嘴上还否认)。我不需要再读;无论口头还是笔下,我都能即兴来一套最纯正的克里希那穆提,一套滴水不漏、环环相扣的话语。可话语再漂亮、再严密,也从来一刻都不像管用,无论对我还是对别人。这样过了好几年,我一点也没做出要从中吸取教训的样子。

随着发现冥想,那套行话一夜之间从我身上脱落,不留痕迹。那时我才懂得话语与认识之间的全部差别。

大老板马上校正了准星:克里希那穆提打入冷宫,冥想捧上台面!自然是悄悄地;如今他得换一套全然不同的手法。时势变了:这小子如今在他腿脚间乱窜,而且有时候眼神挺利。想必小子当时正忙着别处。总之,直到五年之后,某口锅炸开了、小子跑来看出了什么事,大老板的把戏才被识破。

说到底这事过去不算久,才两年多一点:那位“不露相的上师”终于被揭穿——又一副伪装进了冷宫!可怜的老板,他几乎要赤条条地见人了。换个说法:那匹“冥想”马顶替了那匹无名的马(那匹马可千万不能叫“克里希那穆提式”!),如今它带来的押注回报实在少得可笑,尤其比起“数学”马当年那些可观的回报——那时老板还押在它身上。他之所以这么久守着那笔烂注,纯粹是惯性——他已经换过一次注,这本来就不多见,而且当时全靠一记猛击的事件才办到\fnoteP{6}{(42)}{\emph{有益的拔除}\par 所说的那桩“猛击”事件,是1969年底发现:我觉得自己属于的那个机构,有一部分经费来自军事部的款项——这与我基本的公理不相容(而且今天仍然不相容)。这一事件是一整串事件(一个比一个更能说明问题!)中的第一桩,其后果是我离开了高等科学研究所(IHES),并一针引一线地彻底换了圈子、换了投入。\par 在IHES的英雄年代,迪厄多内(Dieudonné)和我是它仅有的成员,也是仅有的为它赢得\textbf{信誉}和听众的人:迪厄多内靠编辑《数学出版物》(Publications Mathématiques)(第一卷早在1959年就出了,即莱昂·莫尚(Léon Motchane)创办IHES的次年);我靠“代数几何讨论班”。最初那些年,IHES的存在极其朝不保夕,经费没有着落(靠几家充当赞助人的公司的慷慨),唯一的场所是巴黎梯也尔基金会(Fondation Thiers)借给的一间屋子(借得明显不情愿),只在我开讨论班的日子用〔IHES最近为建所二十五周年印发的一本小册子(尼科·凯珀(Nico Kuiper)好意寄给我一册)对这些\textbf{艰难}的开端只字不提,也许觉得它们不配去年那场大张旗鼓的庆典的庄严〕。我有点觉得自己和迪厄多内一起,是我所属这个机构的“科学”\textbf{共同创办人},而且我打算在此终老!我最终对IHES产生了强烈的认同,而我的离开(由于同事们的漠然)被体验成一种从另一个“家”被拽离,后来才显出是一场解放。\par 回过头看,我意识到我身上想必早就有一种更新的需要,我说不清从何时起。我离开IHES的前一年,我的精力投入忽然翻了个方向:把前一天还烫手的工作、把我最着迷的问题都撂在一边,转而去搞生物学(是在一位生物学家朋友米尔恰·杜米特雷斯库(Mircea Dumitrescu)的影响下)——这肯定不是单纯的\textbf{巧合}。我当时是以在IHES内部长期投入的架势扑进去的(这与该机构多学科的宗旨相符)。那肯定只是对一种深得多的更新之需的宣泄,而这种更新在IHES那“科学蒸笼”的气氛里无法完成,它是在我前面提到过的那场“觉醒的瀑布”中发生的。前后共有七次,最后一次在1982年。“军费”那段插曲是天意,它触发了这些“觉醒”中的第一次。军事部,和我IHES的老同事们一样,最终都得到了我全部的感激!} (42)。老板们不太喜欢换注——而这一次甚至算是往回走,回到上一笔注。

从1973年我退居乡间起,这匹新马带来的回报,比起从前那匹,就开始少得可怜。三年后冥想的意外降临,又让它们略略活络。甚至还有1979年3月到7月那阵令人头晕的高峰——这一段我在这里不多说——那时我又一次显出使徒的模样,这一回是一种既古老又新鲜的智慧的使徒;这智慧唱在我自己作的一部诗作里,而我最终没有把它交到出版商手上\fnoteP{7}{(43)}{\emph{“我自己作的那部诗作”} 里面有许多我第一手就知道的东西,它们今天在我看来,在我的生命里、以及在一般“生命”里,和它写成时一样重要——当时我是打算发表的。我之所以没发表,主要是因为后来我发觉,它的形式患有一种刻意“作诗”的毛病,以致它的整体构思过于搭建,许多段落缺少自然,有些地方甚至僵硬或浮肿得叫人难受。这种时而浮夸的形式,是我当时心性的映照;在那里,说到底常常是“老板”在领着跳舞——当然是笨重地领……} (43)。可是两年之后,上师彻底歇了业,就好像“冥想”这匹马也断了一条腿(就它带给老板的回报而言)——不管有没有手法,都没法再扮上师了!

% ---------------- p.217 ----------------
\origpage{121} 这以后就没怎么再拖——三条腿的马进了冷宫,连同使徒诗人、“不是上师的上师”,以及“不敢说出自己名字的克里希那穆提”。数学万岁!

下文如何,我们饶有兴致地等着……

% ---------------- p.218 ----------------
% (原书 p.218 为章末空白页)

